% 来源：OEBPS/Text/Section0001_split_002.xhtml#sigil_toc_id_21

\BookChapter{12}
\section[第12章\q 投资营销谎言7：“我恨年金，你也应该恨”]{第12章\\ 投资营销谎言7：“我恨年金，你也应该恨”}

\begin{BookQuote}
美联储主席去年最大的资产是两笔年金。

\hfill ——《美联储主席伯南克的个人理财并不是在装模作样》（《今日美国》2008年7月21日）
\end{BookQuote}

\subsection{爱它还是恨它}

我在网上无意间看到一个在线广告：“我恨年金，你也应该恨。”这是典型的互联网钓鱼式标题，其实它植入的是一篇免费报告，讲的是企业年金这种投资如何糟糕，而采用股票和债券组合是一个好得多的投资策略，因为这种组合长期投资会增值更多，也更安全。当然，广告商已经准备好了要兜售基金产品，其实就是销售它的专业选股技能，它当然要收取一笔管理费。广告标题里没有提到的是，那个花钱打广告的商家是做主动投资的选股机构。那些投资大师，沃伦·巴菲特、约翰·博格、瑞·达利欧、戴维·斯文森都已经给我们讲过了，那些学术研究成果也告诉我们主动管理型基金并不能持续战胜市场。主动管理型基金的业绩还比不上一个简单的指数基金。指数基金的费用要便宜500\%\textasciitilde{}3000\%，业绩却更好。但这种市场营销策略通常很有效，是不是？把你的主动管理型基金业绩和那些大家都觉得非常糟糕的产品的业绩一比，你马上会觉得自己投资的产品没有那么差了。

但是，并不是每个人都“恨”年金。

从另一个角度讲，我非常吃惊地发现，美国联邦储备系统（美联储）前主席本·伯南克在任期间可以说是金融界最有影响力的人，在个人理财计划方面，他非常重视使用年金。伯南克必须先披露个人投资，才能够担任美联储主席。这些公开披露的信息表明，他持有的股票和债券比例相对较低，而他的年金是最大的两笔重仓投资。我马上产生了这样的想法：“难道他知道我并不知道的东西吗？”

那么，他知道的我不知道的东西是什么？

年金是不是自从切片面包发明以来最好的东西？年金只是一种生意，仅对销售年金的保险公司和经纪人有好处吗？答案其实取决于你拥有的年金类型，以及保险公司向你收取的费用。让我们进一步深入研究。

在写本书的过程中，我一直在寻找世界上最受人尊敬的智慧头脑，以探求最好的路径，让读者能够锁定一个有保证的终身收入现金流，一张够你吃上一辈子的支票，从而再也不用为钱去工作。毕竟，这是我们做投资的最大原因。当我挑选要采访的理财大师时，戴维·巴贝尔这个名字一直“持续上升，接近首选地位”。你还记得我在上一章中讲过的吗？他是拥有多个博士学位的沃顿商学院教授，为美国劳工部提供咨询建议，评估过一份关于目标日期基金的研究。

2003年年初，戴维在一份报告里讲他个人的投资理财故事，他是如何拆穿他那些华尔街同行的建议的，那些人鼓励他自己投资，选择股票和债券，追求更高的增值，以此创造一个终身收入计划。相反，他没有冒一分钱的险去投资股票和债券，而是运用一系列有保证收入的年金。年金的长期业绩令人吃惊，让他得到了一份安全且有保障的退休收入。这是他想得到，也值得得到的一份终身收入计划。他运用的年金还能给他100\%的本金安全保证，所以他在2000年和2008年市场暴跌的时候一分钱也没有亏。相反，他现在舒舒服服地享受着退休生活，和老婆、孩子、孙子、孙女开开心心地在一起，内心完全安定平静。因为他知道，这份终身收入计划让他的钱永远也花不完。

我飞到费城，去和巴贝尔博士见面，准备给他做一个小时的访谈，结果访谈变成了四个小时。他的投资策略，我们会在后边重点介绍，它非常强大有力，而且非常简单。这种“内心的宁静”因素确实应该考虑进来，我可以理解。正是他采取的策略能让他享有这种自由。当我结束访谈离开时，我对年金的看法发生了彻底的改变，至少对某些种类的年金的看法是这样的。

他说得非常清楚：“并不是所有的年金都是生而平等的。”年金有很多不同的类型，每一种都有自己独特的好处和缺陷，有些年金你确实应该痛恨，但是把所有年金都归为同一类，不假思索地歧视和痛恨是非常错误的，毕竟年金是唯一经受过2000多年长期考验的金融工具。

\subsection{朱利叶斯·恺撒保险公司}

历史上第一份终身收入年金诞生于2000年前的罗马帝国。罗马市民和战士会把钱存进一个资金池，那些活得最久的人得到的收入会逐年增加，那些非常不幸、没能长寿的人也就得不到钱了。当然，政府会从中抽走一小部分，恺撒的必须归还恺撒。

拉丁语里“annua”这个单词是我们现在的英语单词“annual”的来源，因为最早的罗马人是每年领取收入的。当然，这也是年金（annuity）这个单词的来源。讲这种烦琐的小事和投资有什么关系呢？

17世纪，欧洲政府使用了同样的年金概念（称为联合养老保险），来为战争和公共项目建设筹资。当然，政府要从总金额中抽取一部分费用。在现代世界，这些产品的数学推算和基础仍然完全一样，唯一的改变是信用评级最高的保险公司取代了政府的角色，其中有很多保险公司做年金业务的历史超过了100年。这些保险公司经受了历史的长期考验，包括大萧条、经济衰退、世界大战，以及最近的信贷危机。

但是我们必须小心谨慎，不同类型的年金之间有很大差别。2000多年过去了，年金差不多还是老样子，只有一种版本——像经典可口可乐一样的金融解决方案。它就是一个简单的合同，是你跟保险公司签订的。你把你的钱给保险公司，它向你保证用你的钱，为你创造一份有保证的收入或者收益。你按照合同交完全部费用之后，要决定什么时候开始领取年金。你等的时间越长，能领到的年金就越多。在你购买年金的那一天，你就有一个计划表，上面有明确给付年金的日期和金额，根本不需要你去猜测。

\subsection{这是变化，还是一种进步？}

最近这50年，年金已经演化出了很多不同的类型，完全不同于恺撒时期年金最原始的样子。有时这种演变是一种好事，但有时它最终成了一种突变。

可以肯定地说，糟糕的产品要比好的产品多得多。正如约翰·博格所说的：“从概念上来看，我还是年金的推荐者，但是你最好先仔细地看看所有的细节条款，再决定是否购买。”让我们直接切入主题，你应该避免哪种年金？

\subsection{变额年金都是很糟糕的年金}

2012年，超过1500亿美元的变额年金被卖了出去。为了正确地看待这个问题，我要告诉你，1500亿美元只比苹果公司2012年的销售收入低一点儿。变额年金已经演变成很多大型证券公司和银行最喜欢销售的产品，因为变额年金能给它们带来很多佣金。变额年金到底是什么东西？简言之，它就是一个保险合同，你交的所有保费都会被投资到公募基金（也被称为子账户），是的，还是那种业绩低于市场平均水平但是收费高得离谱的公募基金。但是这一次，投资者所购买的这些公募基金是从一个年金“包”里面买的。为什么会有人愿意通过年金来投资公募基金？因为年金产品有特别的税收优惠，用年金里面的钱投资增值可以递延纳税，就像401(k)账户或者个人退休金账户那样。这种规定特别有吸引力，于是那些推销广告就来了，如果你的401(k)账户或者个人退休金账户缴费都超过了上限，但你还有多余的资本想要投资，那么你可以选择年金。但是你买了年金，不仅要支付过高的费用给包在里面的业绩低下的公募基金，而且年金本身还要另外收取一定的费用。

\subsection{费上加费}

那么年金还有什么吸引力呢？为什么有些人会买那些“包”在年金里面的公募基金，他们只是为了避税吗？大部分变额年金保证，即使账户投资的基金市值下跌，你的受益人还是会至少得到相当于最初投入本金金额的年金。所以，如果你投入了10万美元购买变额年金，但后来年金投资的公募基金的市值跌到了只有2万美元，在你去世之后，你的孩子还是能领到10万美元的年金。听起来它不像是一笔很糟糕的买卖，但是后来你才意识到，其实你是购买了一种价格最昂贵的人寿保险。

在本书的第二部分第7章，我们列了一张基金费用一览表，里面都是你购买一只主动管理型公募基金所需要支付的费用。我们解释过，这些费用会如何大幅地拖累你的投资业绩表现。在这里简单回顾一下，所有这些费用的总额接近于每年3.1\%，这是《福布斯》杂志一篇文章的研究结论。［持有基金用的是可以税收递延的账户，比如401(k)账户、个人退休金账户和变额年金。］

每投资10万美元的基金，每年就要支付3100美元的费用。

但是这还没有完。

如果你买的是变额年金，那么你不但要支付上面所列的基金费用，还要额外支付一些费用给保险公司。首先是“死亡费用”\footnote{\label{fn:pdf-12-1}“死亡费用”包含在某些年金或者保险产品里面，用来补偿保险公司在年金合同下承担的各种风险。}，根据晨星公司的研究，平均每年为1.35\%，再加上年金账户的管理费，每年在0.1\%\textasciitilde{}0.5\%。

让我们把所有这些费用加到一起：

基金平均成本：3.1\%（依据《福布斯》杂志的文章数据）

死亡费用：1.35\%（平均）

管理成本：0.25\%（平均）

合计：4.7\%

也就是说，有时候你投资10万美元的年金，每年就要支付4700美元的费用，你先要支付这些费用，剩下的才是你拿到手的净收益。换句话说，如果基金的投资收益率是4.7\%，你付完费用后，实际上一分钱也得不到。你还指望所有这些附加费用帮你避掉投资收益所要交的税款吗？见鬼去吧，扣掉所有这些费用之后，你根本剩不下什么收益了，如果还剩一点儿的话，你还要自己掏钱交税。

\subsection{陷入困境}

尽管大多数人买这些变额年金都亏了钱，但他们仍然担心、害怕，不敢把钱从年金里拿出来，因为年金有死亡收益保证（保证他们的继承人能得到和原始存入的本金相同的年金）。保险合约上往往有很严格的解约费用条款，所以你要提前解约，保险公司就会向你收取很高的解约费。

那么这个规定有什么例外吗？只有两种年金例外。专家告诉我，到目前为止，当一个人需要提高税收效率的时候，它值得考虑。先锋公司和美国教师退休基金会都提供了低成本的变额年金，里面有一个低成本指数基金名单可供你从中选择投资。它们不收取佣金，所以也没有解约费。你如果想兑现，就可以拿钱走人。

\subsection{不是你爷爷买的那种年金}

在本书第五部分，我们将会清楚地对比传统的收入年金和较新型的年金（固定指数年金），即提供一些最高的和最吸引人的收入，保证其不低于其他任何金融产品，与此同时提供100\%的本金安全保护。等你读完了本书，你心里就会很肯定且非常宁静了，因为你知道每个月你打开信箱时，就会收到一张支票（你不用为了得到这张支票而工作）。我们将会加速你通往财务自由的速度，只要我们能够找到方法消除你获得终身收入所需要交纳的税收就行了。你会问自己要怎么做。

办法是拿出来一部分钱，综合运用罗斯个人退休金账户的力量和终身收入年金的力量。这意味着无论政府把所得税提高到什么程度，你都可以非常放心地确定你拿到的钱都是自己能够花的收入。是的，收入都是合法的、安全的、免税的终身收入。它没有波动的部分，也不用担心股市的变动。

这一章的目的不仅要告诉你需要避免的事情，而且要警告你不要陷入那种市场营销谎言的误导，认为所有年金都是坏的。只有唯一的原因让我不愿意再更详细地解释年金的力量，因为你首先需要明白把你的钱放到哪里了——这就是资产配置。只有明白了资产配置，你才会知道什么时候买年金，以及买什么年金对于你来说是合理的。

\subsection{解决方案}

你如果有了一个年金，不管它是什么类型，找个年金专家帮你评估一下，总是有好处的。

\begin{itemize}
\item 发现你现有年金的优点和缺点。
\item 确定你支付的年金的实际费用是多少。
\item 评估你能得到的保证是否是你可以得到的最高水平。
\item 决定你继续保留现有的年金，还是跳出来“变更”成另外一种不同类型的年金。
\end{itemize}

如果你有一份年金，但是你觉得它不是特别好，那么有一个条款叫作“1035转换”，你只需要做一些简单的文书工作，就可以把现金余额从一家保险公司转到另外一家保险公司，而且不用为此接受税收上的惩罚。但是你必须意识到，你现在的年金可能有“解约费”，如果你没有持有这份年金足够长的时间，就要付解约费。也许，你可以为此推延一下转换，直到解约费变得很低或者根本不再征收。另外，你可能要放弃死亡受益保证金。

继续跟我一起前进，现在只剩下最后也是最终的一个投资营销谎言需要我们去发现了。这个谎言也是投资内行认识得最清楚的：你必须承担非常高的风险，才能获得非常高的收益。

让我们来揭开第8个投资营销谎言……
